\section{Introduction}
\label{sec:intro}

Packet capture is a fundamental tool for debugging, troubleshooting,
and fault analysis of network systems. As link bandwidth and traffic
volume continue to increase, a software capture that forwards every
packet to user space for storage no longer scales~\cite{braun2010capture}. A capture tool must therefore select only
the packets relevant to its purpose as early as possible, ideally
inside the kernel datapath~\cite{papadogiannakis2013scap}. The de-facto language for this purpose has
long been tcpdump/libpcap’s
pcap-filter~\cite{pcapfilter_man,mccanne1993bpf}.

Packet structures whose field locations cannot be determined
statically are common in modern networks. Examples include the inner
headers of GRE, GTP-U, and Geneve tunnels, repeated occurrences of the
same protocol, and variable-length structures such as SRv6 segment
lists and TCP options~\cite{rfc2784,gtp_3gpp,rfc8754,rfc8926,rfc9293}. In these
cases, the location of the target field depends on packet contents and
changes from packet to packet.

Existing packet-filter languages provide limited support for such
structures. pcap-filter expresses conditions over fixed L2–L4 header
fields, but it cannot naturally express nested encapsulation,
repeated protocol occurrences, or variable-length protocol structures.
More expressive languages such as the Wireshark display
filter~\cite{wireshark_dfref} can express conditions over nested
encapsulation and variable-length structures. However, they operate on
parsed packets in user-space memory and cannot execute directly in the
kernel datapath. Existing tools thus leave a gap, and user studies of
packet-analysis tools identify protocol coverage and filter
expressiveness as open problems~\cite{sultana2024survey}.

To execute a filter inside
the Linux kernel, it must also be compiled into eBPF bytecode that the
verifier accepts~\cite{kernel_verifier,ebpf_loops,gershuni2019prevail}.
The verifier requires every packet access to be proven safe. For
fields located behind nested encapsulation or variable-length
structures, the required offset is often known only at run time. As a
result, a naive lowering produces packet accesses whose range the
verifier cannot prove, causing the generated bytecode to be rejected.

Even verifier-aware P4-to-eBPF compilers such as
p4c-xdp~\cite{tu2018p4cxdp} target fixed protocol pipelines rather than
filter expressions over arbitrary encapsulation. Existing systems thus
provide either expressive packet descriptions or verifier-accepted
execution, but not both.

We present \kunai, a verifier-safe packet-filter DSL and compiler for
nested encapsulation and variable-length protocol structures. \kunai
describes packet structure and filter conditions through a
protocol-aware \emph{layer-chain} representation and resolves field
references using protocol definitions written in a subset of
P4~\cite{p4spec}.

The compiler lowers dynamic packet accesses into forms that the
verifier can prove safe. It records the start offset of each parsed
header during packet traversal and lowers every variable-length walk
into a compile-time bounded form. These code-generation patterns
translate field references beyond nested encapsulation and
variable-length structures into verifier-accepted eBPF bytecode. We
use the term \emph{verifier-safe} in the sense of Solleza et
al.~\cite{solleza2025verifiersafe}: a well-formed filter should
compile to bytecode that the verifier accepts.

We evaluated \kunai on filters spanning GTP-U, SRv6, Geneve, and TCP
options. The generated bytecode was accepted by the verifier across
six Linux kernel versions (6.1–7.0) on both the XDP and tc attach points,
matched and rejected packets as intended, and incurred only modest
overhead relative to pcap-filter where equivalent filters could be
expressed. These results showed that such packet filters compiled into
verifier-accepted eBPF programs that matched and rejected packets as
intended.